\documentclass{beamer}

\usepackage[utf8]{inputenc}
\usepackage[T2A]{fontenc}
\usepackage[russian]{babel}
\usepackage{amsmath}
\usepackage{amsfonts}
\usepackage{amssymb}
\usepackage{graphicx}
\usepackage{caption}

\usetheme{Madrid}

\setbeamertemplate{caption}[numbered]

\title[]{Построение СК-оптимальных предикторов\\для финансовых временных рядов}
\author[К.\,А. Артемьев]{Автор: К.\,А. Артемьев}
\institute[]{Руководитель: Е.\,Р. Горяинова}

\begin{document}
	
	\begin{frame}
		\titlepage
	\end{frame}
	
	%\begin{frame}{Цель и задачи работы}
	%	\textbf{Цель:} построение СК-оптимального предиктора доходности трёхгодичных облигаций (DGS3) по её лагам и лагам доходности одногодичных облигаций (DGS1).
	%	
	%	\vspace{0.3cm}
	%	\textbf{Задачи:}
	%	\begin{enumerate}
	%		\item Сравнительный анализ качества оценок параметров линейной регрессии, полученных с помощью:
	%		\begin{itemize}
	%			\item метода наименьших квадратов (МНК),
	%			\item метода наименьших модулей (МНМ),
	%			\item M-оценок Хьюбера и Тьюки,
	%		\end{itemize}
	%		при различных типах аддитивных шумов.
	%		\item Исследование рядов DGS1 ($X$) и DGS3 ($Y$) на коинтеграцию с использованием классических и робастных методов.
	%		\item Исследование причинности по Грейнджеру (G-причинности) с помощью теста Вальда.
	%		\item Построение СК-оптимального предиктора $Y_t$ по $Y_{t-1},\dots,Y_{t-p}, X_{t-1},\dots,X_{t-q}$.
	%	\end{enumerate}
	%\end{frame}
	%
	%\begin{frame}{Оценки параметров линейной регрессии: МНК и МНМ}
	%	Рассмотрим модель $y_i = x_i^T \beta + \varepsilon_i$, $i=1,\dots,n$.
	%	
	%	\begin{block}{Метод наименьших квадратов (МНК, OLS)}
	%		Оценка $\hat{\beta}_{OLS}$ минимизирует сумму квадратов остатков:
	%		\[
	%		\hat{\beta}_{OLS} = \arg\min_{\beta} \sum_{i=1}^n (y_i - x_i^T\beta)^2.
	%		\]
	%		Чувствителен к выбросам (не робастен).
	%	\end{block}
	%	
	%	\begin{block}{Метод наименьших модулей (МНМ, LAD)}
	%		Оценка $\hat{\beta}_{LAD}$ минимизирует сумму абсолютных отклонений:
	%		\[
	%		\hat{\beta}_{LAD} = \arg\min_{\beta} \sum_{i=1}^n |y_i - x_i^T\beta|.
	%		\]
	%		Более робастен к выбросам, чем МНК.
	%	\end{block}
	%\end{frame}
	%
	%\begin{frame}{M-оценки}
	%	M-оценка определяется как
	%	\[
	%	\hat{\beta}_M = \arg\min_{\beta} \sum_{i=1}^n \rho\left(\frac{y_i - x_i^T\beta}{\sigma}\right),
	%	\]
	%	где $\rho(\cdot)$ – функция потерь, $\sigma$ – масштаб.
	%	
	%	\begin{block}{Функция потерь Хьюбера}
	%		\[
	%		\rho_{\text{Хьюбер}}(u) = 
	%		\begin{cases}
	%			\frac{1}{2}u^2, & |u| \le k, \\
	%			k|u| - \frac{1}{2}k^2, & |u| > k,
	%		\end{cases}
	%		\]
	%		с параметром $k$ (обычно $k=1.345$ для 95\% эффективности при нормальном шуме).
	%	\end{block}
	%	
	%	\begin{block}{Функция потерь Тьюки (биквадратная)}
	%		\[
	%		\rho_{\text{Тьюки}}(u) = 
	%		\begin{cases}
	%			\frac{c^2}{6}\left(1 - \left[1 - \left(\frac{u}{c}\right)^2\right]^3\right), & |u| \le c, \\
	%			\frac{c^2}{6}, & |u| > c,
	%		\end{cases}
	%		\]
	%		с параметром $c$ (обычно $c=4.685$).
	%	\end{block}
	%\end{frame}
	
	\begin{frame}{Рассматриваемые данные}
		\begin{figure}[h!]
			\centering\includegraphics[scale=0.5]{../source_data_plot.png}
			\caption{Рассматриваются соответственно процентные доходности однолетних и трёхлетних американских облигаций (обозначаются как DGS1 и DGS3).}
		\end{figure}
	\end{frame}
	
	% Слайд 2: Цель и задачи (компактно)
	\begin{frame}{Цель и задачи работы}
		\small
		\textbf{Цель:} построение СК-оптимального предиктора $Y_t$ по $Y_{t-1},\dots,Y_{t-p}, X_{t-1},\dots,X_{t-q}$, где $X_t$ и $Y_t$ --- временные ряды DGS1 и DGS3 соответственно.
		
		\vspace{0.2cm}
		\textbf{Задачи:}
		\begin{enumerate}
			\item Сравнение качества оценок параметров линейной регрессии, полученных методами:
			\begin{itemize}
				\item МНК,
				\item МНМ,
				\item M-оценками Хьюбера и Тьюки,
			\end{itemize}
			при разных типах аддитивных шумов (вычислительный эксперимент).
			\item Исследование рядов $X_t$ и $Y_t$ на коинтеграцию (классические и робастные методы).
			\item Исследование рядов $X_t$ и $Y_t$ на причинность по Грейнджеру с помощью теста Вальда.
			\item Построение СК-оптимального предиктора.
		\end{enumerate}
	\end{frame}
	
	% Слайд 3: МНК и МНМ (два блока в одной колонке)
	\begin{frame}{Оценки параметров: МНК и МНМ}
		\small
		Модель: $y_i = x_i^T \beta + \varepsilon_i$, $i=1,\dots,n$.
		
		\begin{block}{Метод наименьших квадратов (МНК, OLS)}
			\[
			\hat{\beta}_{OLS} = \arg\min_{\beta} \sum_{i=1}^n (y_i - x_i^T\beta)^2
			\]
			Чувствителен к выбросам.
		\end{block}
		
		\begin{block}{Метод наименьших модулей (МНМ, LAD)}
			\[
			\hat{\beta}_{LAD} = \arg\min_{\beta} \sum_{i=1}^n |y_i - x_i^T\beta|
			\]
			Более робастен к выбросам, чем МНК.
		\end{block}
	\end{frame}
	
	% Слайд 4: M-оценки (функции потерь)
	%\begin{frame}{M-оценки}
	%	\small
	%	M-оценка определяется как
	%	\[
	%	\hat{\beta}_M = \arg\min_{\beta} \sum_{i=1}^n \rho\left(\frac{y_i - x_i^T\beta}{\sigma}\right).
	%	\]
	%	
	%	\begin{columns}[T]
	%		\column{0.48\textwidth}
	%		\textbf{Функция Хьюбера}
	%		\[
	%		\rho_{MH}(u) = 
	%		\begin{cases}
	%			\frac{1}{2}u^2, & |u| \le k, \\
	%			k|u| - \frac{1}{2}k^2, & |u| > k,
	%		\end{cases}
	%		\]
	%		$k=1.345$.
	%		
	%		\column{0.48\textwidth}
	%		\textbf{Функция Тьюки (биквадратная)}
	%		\[
	%		\rho_{MT}(u) = 
	%		\begin{cases}
	%			\frac{c^2}{6}\left(1 - \left[1 - \left(\frac{u}{c}\right)^2\right]^3\right), & |u| \le c, \\
	%			\frac{c^2}{6}, & |u| > c,
	%		\end{cases}
	%		\]
	%		$c=4.685$.
	%	\end{columns}
	%\end{frame}
	
	% Слайд 4: M-оценки (без columns, компактно)
	\begin{frame}{M-оценки}
		\small
		M-оценка определяется как
		\[
		\hat{\beta}_M = \arg\min_{\beta} \sum_{i=1}^n \rho\left(\frac{y_i - x_i^T\beta}{\sigma}\right).
		\]
		
		\textbf{Функция потерь Хьюбера:}
		\[
		\rho_{MH}(u) = 
		\begin{cases}
			\frac{1}{2}u^2, & |u| \le k, \\
			k|u| - \frac{1}{2}k^2, & |u| > k,
		\end{cases}
		\qquad k=1.345.
		\]
		
		\vspace{0.2cm}
		\textbf{Функция потерь Тьюки (биквадратная):}
		\[
		\rho_{MT}(u) = 
		\begin{cases}
			\frac{c^2}{6}\left(1 - \left[1 - \left(\frac{u}{c}\right)^2\right]^3\right), & |u| \le c, \\[4pt]
			\frac{c^2}{6}, & |u| > c,
		\end{cases}
		\qquad c=4.685.
		\]
	\end{frame}
	
	\begin{frame}{Коинтеграция временных рядов}
		Пусть $X_t$ и $Y_t$ – нестационарные ряды $I(1)$. Они называются \textbf{коинтегрированными}, если существует такой параметр $\beta$, что линейная комбинация $E_t = Y_t - \beta X_t$ стационарна.
		
		\begin{block}{Расширенный критерий Дики --- Фуллера (ADF)}
			Для проверки стационарности остатков $E_t$ оценивается модель:
			\[
			\Delta E_t = \gamma E_{t-1} + \sum_{i=1}^{p} \theta_i \Delta E_{t-i} + \varepsilon_t.
			\]
			Гипотеза $H_0$: $\gamma = 0$ (наличие единичного корня, нестационарность). Если статистика $\widehat{DF}_0(\hat\gamma)$ меньше критического значения $DF_0$, то $H_0$ отвергается --- ряды коинтегрированы.
		\end{block}
	\end{frame}
	
	\begin{frame}{Результаты исследования на коинтеграцию}
		\begin{itemize}
			\item Построены регрессии $Y_t = \beta_0 + \beta_1 X_t + \varepsilon_t$ классическими и робастными методами (МНК, МНМ, M-оценки Хьюбера и Тьюки).
			\item К остаткам применён расширенный критерий Дики --- Фуллера.
			\item \textbf{Результат:} во всех случаях $p$-значение ADF-теста < 0.05, что свидетельствует о стационарности остатков. Ряды DGS1 и DGS3 коинтегрированы.
			%\item \textbf{Вывод:} ряды DGS1 и DGS3 коинтегрированы – существует долгосрочная равновесная связь между доходностями.
		\end{itemize}
	\end{frame}
	
	\begin{frame}{Причинность по Грейнджеру (G-причинность)}
		Временной ряд $X_t$ является причиной по Грейнджеру для $Y_t$, если прошлые значения $X_t$ помогают предсказывать $Y_t$ при учёте собственных лагов $Y_t$, т.\,е.
		\begin{equation*}
			\mathbb{E} \left[ Y_t \mid Y_{\tau}, \, \tau<t \right] \ne \mathbb{E} \left[ Y_t \mid Y_{\tau}, \, X_{\tau}, \, \tau<t \right].
		\end{equation*}
		
		%Рассмотрим VAR-модель порядка $p$:
		Будем рассматривать следующую модель:
		\[
		Y_t = \alpha_0 + \sum_{i=1}^p \alpha_i Y_{t-i} + \sum_{i=1}^q \alpha_{p+i} X_{t-i} + \varepsilon_t.
		\]
		
		\textbf{Нулевая гипотеза} об отсутствии причинности $X \not\to Y$:
		\[
		H_0: \alpha_{p+1} = \alpha_{p+2} = \ldots = \alpha_{p+q} = 0.
		\]
	\end{frame}
	
	\begin{frame}{Тест Вальда для проверки причинности}
		Статистика Вальда для линейных ограничений $R \alpha = r$:
		\[
		\widehat{W}(\hat{\alpha}) = (R\hat{\alpha} - r)^T \left[ R \widehat{\text{Cov}}(\hat{\alpha}) R^T \right]^{-1} (R\hat{\alpha} - r).
		\]
		В нашем случае $R$ --- матрица с размерностью $q \times (p+q+1)$, выделяющая коэффициенты при $X_{t-i}$, а $r$ --- столбец нулей.
		
		При $H_0$ статистика $W$ асимптотически распределена как $\chi^2_q$.
		
		\vspace{0.3cm}
		\textbf{Правило принятия решения:}
		Если $\widehat{W}(\hat{\alpha}) > \chi^2_q(1-a)$ (или p-value~$< a$), гипотеза об отсутствии причинности отвергается.
		
		\textbf{Результаты:}
		\begin{itemize}
			\item $X$ (доходность однолетних облигаций) является G-причиной для $Y$ (доходности трёхлетних облигаций).
			\item Обратная причинность не обнаружена.
		\end{itemize}
	\end{frame}
	
	%\begin{frame}{Результаты теста на причинность по Грейнджеру}
	%	Построена VAR-модель с лагом $p$, выбранным по критерию Акаике (AIC).
	%	
	%	\begin{table}[h]
	%		\centering
	%		\begin{tabular}{|c|c|c|}
	%			\hline
	%			Гипотеза & Статистика Вальда & p-value \\
	%			\hline
	%			$X \to Y$ & 15.23 & 0.001 \\
	%			$Y \to X$ & 2.14 & 0.544 \\
	%			\hline
	%		\end{tabular}
	%		\caption{Результаты тестирования (примерные данные)}
	%	\end{table}
	%	
	%	\vspace{0.3cm}
	%	\textbf{Результаты:}
	%	\begin{itemize}
	%		\item $X$ (доходность 1-летних облигаций) является G-причиной для $Y$ (доходности 3-летних облигаций).
	%		\item Обратная причинность не обнаружена.
	%	\end{itemize}
	%\end{frame}
	
	\begin{frame}{Построение СК-оптимального предиктора}
		\textbf{Среднеквадратически оптимальный предиктор} – условное математическое ожидание:
		\[
		\hat{Y}_t = \mathbb{E}[Y_t \mid \mathcal{F}_{t-1}],
		\]
		где $\mathcal{F}_{t-1}$ – информация о прошлых значениях $Y$ и $X$.
		
		В предположении линейности:
		\[
		\hat{Y}_{t+h} = \hat{\alpha}_0 + \sum_{i=0}^{p-1} \hat{\alpha}_i Y_{t-i} + \sum_{i=0}^{q-1} \hat{\alpha}_{p+i} X_{t-i}.
		\]
		
		Параметры оцениваются с помощью МНК или робастных методов (например, M-оценок) для обеспечения устойчивости к выбросам.
		
		На основе коинтеграции и G-причинности выбирается спецификация: $p=2$, $q=2$ (по AIC). Предиктор демонстрирует высокую точность на тестовой выборке.
	\end{frame}
	
	\begin{frame}{Заключение}
		\begin{itemize}
			\item В ходе вычислительного эксперимента показано, что робастные методы (МНМ, M-оценки) превосходят МНК при наличии тяжёлых хвостов и выбросов.
			\item Подтверждена коинтеграция временных рядов DGS1 и DGS3.
			\item Выявлена односторонняя причинность по Грейнджеру:\\DGS1 $\to$ DGS3.
			%\item Построен СК-оптимальный линейный предиктор $Y_t$, учитывающий как собственные лаги, так и лаги $X_t$.
			%\item Полученные результаты могут быть использованы для прогнозирования доходности средне- и долгосрочных облигаций.
		\end{itemize}
	\end{frame}
	
\end{document}